% block: condition-number-heading page:17 type:section
\subsection{{条件数}}\label{sec:condition-number}
% block: perturbation-equation page:17 type:ordinary-paragraph
{考虑 $A\bm x=\bm b$; 如果矩阵 $A$ 和右端向量 $\bm b$ 发生微小变化, 解会改变多少? 为回答这一问题, 设 $\Delta A$ 和 $\Delta\bm b$ 分别为 $A$ 和 $\bm b$ 的扰动. 以 $\Delta\bm x$ 表示的解的变化满足下列关系:}
\[
(A+\Delta A)(\bm x+\Delta\bm x)=\bm b+\Delta\bm b.
\]
{利用 $A\bm x=\bm b$ 得到}
\[
\Delta\bm x=A^{-1}[-\Delta A(\bm x+\Delta\bm x)+\Delta\bm b].
\]
{利用 Theorem~\ref{thm:matrix-norm-properties}-(\ref{part:matrix-norm-action}) 和三角不等式 (\ref{cond:vector-norm-triangle}), 可由下式估计 $\bm x$ 的 $p$-范数:}
\[
\|\Delta\bm x\|_p\le\|A^{-1}\|_p(\|\Delta A\|_p\|\bm x+\Delta\bm x\|_p+\|\Delta\bm b\|_p).
\]
{因此, 有}
\begin{equation}\label{eq:relative-perturbation-bound}
\frac{\|\Delta\bm x\|_p}{\|\bm x+\Delta\bm x\|_p}
\le\|A\|_p\|A^{-1}\|_p
\left(\frac{\|\Delta A\|_p}{\|A\|_p}+
\frac{\|\Delta\bm b\|_p}{\|A\|_p\|\bm x+\Delta\bm x\|_p}\right).
\end{equation}
% block: condition-number-definition page:17 type:ordinary-paragraph
{这里,}
\begin{equation}\label{eq:p-norm-condition-number}
\kappa_p(A):=\|A\|_p\|A^{-1}\|_p
\end{equation}
{称为 $p$-范数条件数. 不等式 \eqref{eq:relative-perturbation-bound} 表明, 如果条件数很大, 计算得到的近似解可能远离精确解. 注意, 这一事实与求解线性方程组所使用的算法无关.}
% block: ill-conditioning page:17 type:ordinary-paragraph
{如果条件数很大, 则称该线性方程组是病态的. 我们将在 Chapter~\ref{chap:krylov-classification} 看到, 条件数也会影响 Krylov 子空间方法的收敛速度.}
{本书经常使用 2-范数条件数 $\kappa_2(A)$, 并将其简记为 $\kappa$ 或 $\operatorname{cond}(A)$.}
